\documentclass[11pt,a4paper]{article}

\usepackage{geometry}
\geometry{margin=2.5cm}
\usepackage{fontspec}
\setmainfont{Times New Roman}
\usepackage{amsmath,amssymb}
\usepackage{booktabs}
\usepackage{tabularx}
\usepackage{graphicx}
\usepackage{xcolor}
\usepackage{listings}
\usepackage{cite}
\usepackage{fancyhdr}
\usepackage{microtype}
\usepackage{caption}
\usepackage{subcaption}
\usepackage{url}
\def\UrlBreaks{\do\/\do-\do\_}
\usepackage{hyperref}

% Headheight configuration
\setlength{\headheight}{14.5pt}
\pagestyle{fancy}
\fancyhf{}
\rhead{\footnotesize Nghiên Cứu Hệ Thống Phenikaa Calendar Exporter (Web3)}
\lhead{\footnotesize Báo Cáo Nghiên Cứu Học Thuật}
\cfoot{\thepage}

% Color palette
\definecolor{primaryblue}{RGB}{18, 52, 102}
\definecolor{codebg}{RGB}{245, 247, 250}
\definecolor{codeframe}{RGB}{209, 213, 219}
\definecolor{codegreen}{RGB}{16, 124, 65}
\definecolor{codegray}{RGB}{107, 114, 128}

\hypersetup{
    colorlinks=true,
    linkcolor=primaryblue,
    citecolor=primaryblue,
    urlcolor=primaryblue,
    pdftitle={Nghien Cuu Kien Truc Phi Tap Trung Web3 Cho He Thong Quan Ly Lich Hoc Tap},
    pdfauthor={Nguyen Minh Thai, Huynh Quang Thinh}
}

% Table columns
\newcolumntype{Y}{>{\raggedright\arraybackslash}X}
\newcolumntype{L}{>{\raggedright\arraybackslash}X}
\newcolumntype{C}{>{\centering\arraybackslash}X}

% Custom listings for JSON and Python
\lstdefinelanguage{json}{
    basicstyle=\ttfamily\footnotesize,
    numbers=left,
    numberstyle=\tiny\color{codegray},
    stepnumber=1,
    numbersep=6pt,
    showstringspaces=false,
    breaklines=true,
    frame=single,
    backgroundcolor=\color{codebg},
    rulecolor=\color{codeframe},
    stringstyle=\color{primaryblue},
    keywordstyle=\color{codegreen}
}

\lstdefinelanguage{solidity}{
    basicstyle=\ttfamily\footnotesize,
    numbers=left,
    numberstyle=\tiny\color{codegray},
    stepnumber=1,
    numbersep=6pt,
    showstringspaces=false,
    breaklines=true,
    frame=single,
    backgroundcolor=\color{codebg},
    rulecolor=\color{codeframe},
    keywordstyle=\color{primaryblue}\bfseries,
    commentstyle=\color{codegray}
}

% Vietnamese local terms
\renewcommand{\abstractname}{Tóm Tắt}
\renewcommand{\refname}{Tài Liệu Tham Khảo}
\renewcommand{\figurename}{Hình}
\renewcommand{\tablename}{Bảng}

\title{\LARGE\textbf{NGHIÊN CỨU VÀ HIỆN THỰC HÓA KIẾN TRÚC PHI TẬP TRUNG (WEB3)\\CHO HỆ THỐNG QUẢN LÝ VÀ ĐỒNG BỘ LỊCH HỌC TẬP}\\
\vspace{0.4cm}
\large\textbf{Nghiên cứu điển hình trên Hệ thống Phenikaa Calendar Exporter}}

\author{
    \textbf{Nguyễn Minh Thái}\textsuperscript{1}\quad\quad\textbf{Huỳnh Quang Thịnh}\textsuperscript{2} \\
    Khoa Khoa học Máy tính, Trường Công nghệ Thông tin, Đại học Phenikaa \\
    Hà Nội, Việt Nam \\
    Email: \textsuperscript{1}\href{mailto:b4iterdev@gmail.com}{\texttt{b4iterdev@gmail.com}},\quad
    \textsuperscript{2}\href{mailto:thinhbioaxetic06@gmail.com}{\texttt{thinhbioaxetic06@gmail.com}} \\
    Mã nguồn: \href{https://github.com/b4iterdev/PhenikaaCalendarExporter-Web3}{\texttt{github.com/b4iterdev/PhenikaaCalendarExporter-Web3}}
}

\date{\today}

\begin{document}

\maketitle

\begin{abstract}
Quản lý lịch học và lịch thi là nhu cầu thiết yếu đối với sinh viên trong các cơ sở giáo dục đại học. Tuy nhiên, các hệ sinh thái phần mềm hỗ trợ đào tạo hiện nay phần lớn vận hành dựa trên các dịch vụ định danh tập trung (Identity Providers - IdP) và hệ thống lưu trữ máy chủ cục bộ. Cấu trúc này tiềm ẩn các rủi ro về quyền riêng tư dữ liệu, nguy cơ phụ thuộc nền tảng và rủi ro điểm lỗi đơn lẻ (Single Point of Failure).

Báo cáo nghiên cứu học thuật này đề xuất và hiện thực hóa một kiến trúc hệ thống phi tập trung (\textit{Decentralized Architecture}) ứng dụng công nghệ Web3 vào bài toán trích xuất, lưu trữ và điều phối truy cập dữ liệu thời khóa biểu đại học thông qua hệ thống \textbf{Phenikaa Calendar Exporter (Web3 Version)}. Đóng góp học thuật và giải pháp kỹ thuật của nghiên cứu bao gồm:
(1) Mô hình định danh tự chủ (Self-Sovereign Identity) dựa trên giao thức \textit{Sign-In with Ethereum} (SIWE - EIP-4361), kết hợp quy tắc chuẩn hóa địa chỉ EIP-55 Checksum bảo đảm tính toàn vẹn xác thực mật mã;
(2) Cơ chế lưu trữ phân tán nội dung bất biến trên mạng IPFS (InterPlanetary File System) định danh qua CIDv1, hỗ trợ cơ chế bộ đệm xác định phục vụ môi trường kiểm thử độc lập;
(3) Mô hình điều phối truy cập và phân bổ hạn mức API tầng ứng dụng dựa trên hợp đồng thông minh on-chain (Token-Gating) thông qua hàm truy vấn số dư chuẩn ERC-20 / ERC-721 trên mạng EVM Layer-2;
(4) Thiết kế và kiểm chứng thực nghiệm mô hình vận hành liên tục không gián đoạn dịch vụ (\textit{RollingUpdate Zero-Downtime}) trên cụm điều phối K3s Kubernetes kết hợp luồng tự động hóa GitOps Flux v2.

Hệ thống được đánh giá thực nghiệm toàn diện thông qua 158 kịch bản kiểm thử tự động (đạt tỷ lệ vượt qua 100\%) và vận hành ổn định trên môi trường thực tế tại tên miền \url{https://pca-web3.b4iterdev.net.eu.org}. Kết quả chứng minh kiến trúc đề xuất giải quyết hiệu quả các hạn chế của mô hình tập trung, bảo đảm quyền tự chủ dữ liệu và duy trì tính liên tục của dịch vụ.
\end{abstract}

\vspace{0.5cm}
\noindent\textbf{Từ khóa:} Web3, Sign-In with Ethereum (SIWE), EIP-4361, IPFS, Content Identifier (CID), Token Gating, Kubernetes (K3s), Zero-Downtime Deployment, RollingUpdate, GitOps Flux v2.

\newpage
\tableofcontents
\newpage

\section{Giới Thiệu \& Tổng Quan Vấn Đề Nghiên Cứu}

\subsection{Bối cảnh thực tiễn}
Trong các cơ sở giáo dục đại học, lịch thời khóa biểu và lịch thi đóng vai trò trung tâm trong kế hoạch học tập của sinh viên. Tại trường Đại học Phenikaa, hệ thống Quản lý đào tạo (QLĐT) cung cấp lịch học thông qua giao diện web ASP.NET. Để sử dụng lịch trên các thiết bị cá nhân hoặc tích hợp vào các ứng dụng quản lý thời gian như Google Calendar, Apple Calendar và Microsoft Outlook, sinh viên cần công cụ tự động trích xuất và chuyển đổi dữ liệu.

Hệ thống \texttt{PhenikaaCalendarExporter} được phát triển nhằm tự động hóa quy trình phân tích dữ liệu phiên làm việc, truy vấn API nội bộ của cổng thông tin và chuyển đổi thành các định dạng tiêu chuẩn gồm iCalendar (\texttt{.ics}), Excel (\texttt{.xlsx}) và JSON (\texttt{.json}).

\subsection{Hạn chế của mô hình kiến trúc tập trung}
Mô hình vận hành của các ứng dụng tích hợp giáo dục truyền thống thường đối mặt với ba rào cản chính:
\begin{enumerate}
    \item \textbf{Phụ thuộc định danh tập trung (Centralized Identity Dependency):} Việc phụ thuộc vào các nhà cung cấp bên thứ ba (Google OAuth, OIDC) buộc người dùng phải cung cấp thông tin định danh cá nhân (email, hồ sơ cá nhân). Khi tài khoản bên thứ ba gặp sự cố khóa hoặc giới hạn chính sách, người dùng mất hoàn toàn quyền truy cập dịch vụ.
    \item \textbf{Rủi ro dữ liệu lưu trữ máy chủ cục bộ (Single Server Storage Risk):} Các tệp dữ liệu lịch thường được lưu trên ổ đĩa cục bộ của máy chủ dịch vụ. Mô hình này thiếu tính bất biến, không hỗ trợ xác minh tính toàn vẹn độc lập và đối mặt nguy cơ mất mát dữ liệu khi hạ tầng lưu trữ vật lý gặp sự cố.
    \item \textbf{Kiểm soát truy cập thiếu tính minh bạch:} Việc phân quyền và giới hạn lưu lượng (rate limiting) được quản lý kín trong cơ sở dữ liệu nội bộ máy chủ, người dùng không thể tự xác minh hay chủ động chuyển giao quyền hạn truy cập theo chuẩn mở.
\end{enumerate}

\subsection{Mục tiêu và đóng góp của nghiên cứu}
Nghiên cứu này tập trung thiết kế và hiện thực hóa kiến trúc Web3 cho hệ thống trích xuất lịch học, mang lại các đóng góp chính:
\begin{itemize}
    \item \textbf{Mô hình định danh tự chủ:} Triển khai giao thức EIP-4361 (SIWE), cho phép người dùng xác thực thông qua chữ ký mật mã khóa công khai EVM mà không chia sẻ thông tin nhận dạng cá nhân.
    \item \textbf{Hạ tầng lưu trữ phi tập trung:} Đưa các tệp xuất dữ liệu lịch lên hệ thống tệp phân tán IPFS với mã định danh nội dung bất biến CIDv1, hỗ trợ truy xuất phân tán ngang hàng.
    \item \textbf{Cơ chế Token-Gating tầng ứng dụng:} Ứng dụng hợp đồng thông minh chuẩn ERC-20 / ERC-721 trên mạng EVM Layer-2 để phân quyền và điều phối hạn mức gọi API v1.
    \item \textbf{Mô hình vận hành liên tục Zero-Downtime:} Hiện thực hóa quy trình triển khai phân tán không gián đoạn (\texttt{RollingUpdate}) trên cụm K3s Kubernetes kết hợp GitOps Flux v2.
\end{itemize}

\section{Thiết Kế Kiến Trúc Hệ Thống Phi Tập Trung}

Kiến trúc tổng thể của hệ thống Web3 được phân tách thành 4 lớp cốt lõi: Lớp Định danh (Identity Layer), Lớp Lưu trữ Phân tán (Storage Layer), Lớp Điều phối \& Hợp đồng Thông minh (Smart Contract \& Access Layer), và Lớp Hạ tầng Điều phối (Infrastructure Layer). Bảng~\ref{tab:architecture_layers} tổng hợp các thành phần kỹ thuật chính của kiến trúc này.

\begin{table}[htbp]
\centering
\caption{Các phân tầng kỹ thuật trong Kiến trúc Phi tập trung của Hệ thống}
\label{tab:architecture_layers}
\small
\begin{tabularx}{\textwidth}{p{4.2cm} >{\raggedright\arraybackslash}X >{\raggedright\arraybackslash}X}
\toprule
\textbf{Phân tầng Kiến trúc} & \textbf{Công nghệ \& Chuẩn triển khai} & \textbf{Vai trò chức năng} \\
\midrule
\textbf{Lớp Định danh (Identity)} & SIWE EIP-4361, EIP-55 Checksum, ECDSA & Xác thực phi tập trung bằng cặp khóa bí mật/công khai, không lưu PII. \\
\textbf{Lớp Lưu trữ (Storage)} & IPFS, CIDv1, Pinata Distributed Pinning & Lưu trữ tệp lịch học bất biến, định danh nội dung phân tán toàn cầu. \\
\textbf{Lớp Điều phối (Access)} & Smart Contracts (ERC-20/721), JSON-RPC & Kiểm soát hạn mức gọi API tầng ứng dụng dựa trên số dư on-chain. \\
\textbf{Cổng Lập trình (API Portal)} & RESTful API v1, Granular Scopes, HMAC & Cung cấp giao diện tương tác lập trình và cấp phát token an toàn. \\
\textbf{Lớp Vận hành (Infra)} & K3s Kubernetes, GitOps Flux v2, Docker & Điều phối cụm container với độ sẵn sàng cao, cập nhật không downtime. \\
\bottomrule
\end{tabularx}
\end{table}

\subsection{Định danh Tự chủ: Sign-In with Ethereum (SIWE / EIP-4361)}
Thay thế quy trình chuyển hướng IdP truyền thống, giao thức SIWE \cite{eip4361} sử dụng cặp khóa mật mã để xác thực danh tính người dùng một cách an toàn và bảo mật.

\subsubsection{Giao thức bắt tay mật mã (Handshake Protocol)}
Quy trình xác thực được thiết kế theo các bước:
\begin{enumerate}
    \item \textbf{Cấp phát Nonce mật mã:} Trình duyệt gửi yêu cầu \texttt{GET /auth/siwe/nonce}. Máy chủ sinh ngẫu nhiên chuỗi hex an toàn (\texttt{secrets.token\_hex(16)}), đóng gói vào cookie \path{phenikaa_siwe_nonce} ký mã hóa HMAC với TTL 300 giây.
    \item \textbf{Cấu trúc thông điệp EIP-4361:} Ứng dụng client truy vấn địa chỉ ví và mã mạng từ provider EIP-1193, khởi tạo thông điệp chuẩn hóa chứa các trường: miền dịch vụ (\texttt{domain}), địa chỉ ví (\texttt{address}), chuỗi ngẫu nhiên (\texttt{nonce}), thời điểm khởi tạo (\texttt{issuedAt}) và Chain ID:
    \begin{quote}
    \small\ttfamily
    \{domain\} wants you to sign in with your Ethereum account:\\
    \{address\}\\
    \\
    Sign in with Ethereum to Phenikaa Calendar Web3.\\
    \\
    URI: \{origin\}\\
    Version: 1\\
    Chain ID: \{chainId\}\\
    Nonce: \{nonce\}\\
    Issued At: \{issuedAt\}
    \end{quote}
    \item \textbf{Ký ngoài chuỗi (Off-chain Cryptographic Signature):} Người dùng thực hiện ký chuỗi thông điệp bằng khóa riêng thông qua thuật toán ECDSA (phương thức \texttt{personal\_sign}). Khóa riêng tư được bảo vệ tuyệt đối trong ví phần cứng hoặc tiện ích ví cá nhân.
    \item \textbf{Xác minh chữ ký và Khởi tạo phiên:} Client gửi payload gồm \texttt{\{message, signature\}} đến máy chủ tại endpoint \path{/auth/siwe/verify}. Máy chủ khôi phục địa chỉ công khai từ chữ ký, đối soát tính hợp lệ của Nonce và Domain, sau đó cấp phát cookie phiên bảo mật gắn với định danh ví \path{eth:<address>}.
\end{enumerate}

\subsubsection{Chuẩn hóa địa chỉ EIP-55 Checksum}
Để đảm bảo tính nhất quán giữa dữ liệu địa chỉ đầu vào và chữ ký mật mã, hệ thống áp dụng cơ chế chuẩn hóa EIP-55 \cite{eip55}:
\begin{itemize}
    \item Tách riêng bản văn ký gốc (\path{raw_message}) làm đầu vào phục hồi chữ ký thông qua thuật toán \texttt{recover\_message}.
    \item Tự động chuẩn hóa địa chỉ ví sang định dạng checksum theo hàm băm Keccak-256 (\path{to_checksum_address}) trước khi tiến hành xác minh ngữ pháp.
\end{itemize}

\subsection{Lưu Trữ Phân Tán IPFS}
Hệ thống tích hợp lớp dịch vụ \texttt{IpfsService} cho phép đưa các tệp xuất dữ liệu lịch học lên mạng phân tán IPFS \cite{benet2014ipfs}.

\subsubsection{Cơ chế ghim dữ liệu phân tán (IPFS Pinning)}
Dịch vụ giao tiếp với endpoint cổng phân tán thông qua giao thức HTTP POST đa phần:
\begin{center}
\footnotesize\url{https://api.pinata.cloud/pinning/pinFileToIPFS}
\end{center}
\begin{itemize}
    \item Các tệp lịch học (\texttt{calendar.json}, \texttt{calendar.ics}, \texttt{calendar.xlsx}) được chuyển thành mảng byte và truyền tải tới mạng lưới lưu trữ IPFS.
    \item Dịch vụ ghim phân tán tệp đến các node mạng toàn cầu và phản hồi mã định danh nội dung \texttt{IpfsHash} (CIDv1).
    \item Dữ liệu lịch học có thể truy xuất công khai thông qua cấu trúc URI chuẩn \path{ipfs://<CID>} hoặc qua các cổng IPFS Gateway toàn cầu.
\end{itemize}

\subsubsection{Bộ mô phỏng tính toán CID xác định (Deterministic CID Computation)}
Nhằm phục vụ môi trường kiểm thử tự động độc lập và duy trì tính sẵn sàng cao khi hoạt động ngoại tuyến, \texttt{IpfsService} tích hợp thuật toán tính mã băm xác định từ nội dung tệp:
\begin{equation}
\text{CID}_{\text{mock}} = \text{``bafybeimock''} + \text{SHA-256}(\text{file\_bytes})[0:32]
\end{equation}
Thuật toán này bảo đảm các phép kiểm thử tính toàn vẹn dữ liệu được thực thi với độ trễ tối thiểu mà vẫn duy trì định dạng cấu trúc chuẩn của CIDv1.

\subsection{Cơ chế Token-Gating \& Điều Phối Truy Cập On-chain}
Module \texttt{TokenGatingService} thiết lập cơ chế kiểm soát truy cập tầng ứng dụng dựa trên trạng thái tài sản số on-chain.

\subsubsection{Tương tác Smart Contract qua EVM JSON-RPC}
Hệ thống kết nối trực tiếp đến các node RPC trên mạng EVM Layer-2 (mặc định: Polygon Mainnet \url{https://polygon-rpc.com}, Chain ID: 137). Để tối ưu băng thông và giảm thiểu độ trễ truy vấn mạng, hệ thống sử dụng giao diện nhị phân (ABI) tối giản của hàm truy vấn số dư:
\begin{lstlisting}[language=solidity]
function balanceOf(address _owner) external view returns (uint256 balance);
\end{lstlisting}

\subsubsection{Quy trình xác thực quyền Token-Gate}
Khi một yêu cầu API v1 mang Bearer Token được gửi đến, máy chủ thực hiện chu trình kiểm tra:
\begin{enumerate}
    \item Xác thực mã token trong cơ sở dữ liệu hệ thống thông qua mã băm SHA-256 (\texttt{token\_hash}).
    \item Truy xuất thông tin chủ sở hữu tài khoản gắn với địa chỉ ví \path{eth:0x...}.
    \item Thực hiện truy vấn trực tiếp hàm \texttt{balanceOf(wallet\_address)} trên hợp đồng thông minh đã cấu hình.
    \item Nếu số dư tài sản on-chain thỏa mãn hạn mức tối thiểu (\texttt{min\_balance}), yêu cầu được cấp phép xử lý; ngược lại, hệ thống trả về mã từ chối truy cập HTTP 403 Forbidden kèm mã lỗi chuẩn hóa.
\end{enumerate}

\section{Cổng Phát Triển Tích Hợp \& Giao Diện Trực Tiếp}

Hệ thống cung cấp giao diện quản trị phát triển tại đường dẫn \path{/developers}, phục vụ việc cấp phát khóa truy cập và kiểm thử tương tác API thời gian thực.

\subsection{Cơ chế cấp quyền hạt mịn (Granular Scopes)}
Người dùng có thể tùy biến giới hạn phạm vi tác vụ cho từng API key độc lập:
\begin{itemize}
    \item Quyền đọc thông tin định danh: \texttt{me:read}.
    \item Quyền quản lý token truy cập: \texttt{tokens:read}, \texttt{tokens:write}.
    \item Quyền quản lý phiên và truy xuất sự kiện lịch: \texttt{sessions:read}, \texttt{events:read}.
    \item Quyền kích hoạt tiến trình đồng bộ dữ liệu: \texttt{sync:write}.
\end{itemize}

\subsection{Bảo mật khóa truy cập (One-time Key Reveal)}
Nhằm bảo vệ an toàn tối đa cho người dùng, khóa bí mật thô (\texttt{phk\_...}) chỉ được hiển thị đúng một lần duy nhất tại thời điểm tạo mới. Hệ thống áp dụng nguyên tắc lưu trữ một chiều: máy chủ chỉ ghi nhận mã băm SHA-256 của token và tiền tố 12 ký tự (\texttt{token\_prefix}) phục vụ việc nhận diện trong danh sách quản trị.

\subsection{Bàn thử nghiệm API trực tiếp (Interactive API Playground)}
Bàn thử nghiệm tích hợp hỗ trợ đầy đủ 11 endpoint RESTful của chuẩn API v1, cho phép các nhà nghiên cứu và phát triển:
\begin{itemize}
    \item Đo đạc thời gian trễ phản hồi (latency) bằng mili-giây qua API trình duyệt \path{performance.now()}.
    \item Sinh tự động mã nguồn mẫu thực thi bằng lệnh \texttt{cURL} và JavaScript \texttt{fetch()} tương ứng với từng cấu hình tham số.
    \item Kiểm thử các kịch bản xuất lịch trực tiếp mà không yêu cầu lưu trữ dữ liệu người dùng lâu dài.
\end{itemize}

\section{Hạ Tầng Điều Phối K3s \& Triển Khai Không Gián Đoạn}

Dịch vụ được đóng gói và vận hành trên cụm Kubernetes rút gọn K3s \cite{k3s} kết hợp hệ thống GitOps Flux v2 \cite{fluxcd}.

\subsection{Chiến lược RollingUpdate Zero-Downtime}
Để bảo đảm tính liên tục của các kết nối HTTP/2 dài và ngăn ngừa gián đoạn truy cập trong quá trình cập nhật phiên bản, cấu hình Deployment của hệ thống áp dụng chiến lược \texttt{RollingUpdate}:
\begin{lstlisting}[language=json]
{
  "strategy": {
    "type": "RollingUpdate",
    "rollingUpdate": {
      "maxSurge": 1,
      "maxUnavailable": 0
    }
  },
  "terminationGracePeriodSeconds": 10
}
\end{lstlisting}

\subsubsection{Cơ chế vận hành của chiến lược RollingUpdate}
\begin{itemize}
    \item \textbf{\texttt{maxSurge: 1}:} Cho phép Kubernetes khởi tạo một pod mới bổ sung song song với số lượng replica hiện hành. Pod mới thực hiện kéo container image và kích hoạt ứng dụng.
    \item \textbf{\texttt{maxUnavailable: 0}:} Tuyệt đối không cho phép pod đang chạy bị tắt trước khi pod mới hoàn tất khởi động. Toàn bộ lưu lượng người dùng được duy trì liên tục qua pod cũ.
    \item \textbf{Đầu dò kiểm tra sẵn sàng (\texttt{readinessProbe}):} Ingress Controller chỉ bắt đầu điều phối lưu lượng sang pod mới khi pod vượt qua bài kiểm tra sức khỏe \texttt{GET /healthz} với mã phản hồi HTTP 200.
    \item \textbf{Chấm dứt mềm (Graceful Termination):} Sau khi lưu lượng được chuyển giao sang pod mới, pod cũ nhận tín hiệu \texttt{SIGTERM} và duy trì hoạt động trong 10 giây để hoàn tất các luồng xử lý xuất lịch dang dở.
\end{itemize}

\subsection{Tự động hóa CI/CD và GitOps}
Quy trình xây dựng và phát hành phần mềm được tự động hóa hoàn toàn:
\begin{enumerate}
    \item \textbf{Khâu kiểm thử và đóng gói:} Quy trình GitHub Actions tự động kiểm thử toàn bộ mã nguồn trên môi trường Python 3.12, đóng gói container đa kiến trúc (\texttt{linux/amd64}, \texttt{linux/arm64}) và đẩy lên GitHub Container Registry.
    \item \textbf{Khâu cập nhật cấu hình:} Mã băm bất biến của container image (\path{IMAGE_DIGEST}) được ghi tự động vào kho lưu trữ cấu hình GitOps.
    \item \textbf{Hòa giải trạng thái GitOps:} Bộ điều khiển Flux v2 định kỳ phát hiện thay đổi và đồng bộ trạng thái mong muốn lên cụm K3s mà không cần can thiệp thủ công.
\end{enumerate}

\section{Kết Quả Thực Nghiệm \& Đánh Giá}

\subsection{Đánh giá kiểm thử chức năng tự động}
Hệ thống được kiểm thử tự động với bộ kiểm thử 158 kịch bản bao phủ toàn bộ các module từ phân tích HTML, xử lý phiên, xác thực SIWE, lưu trữ IPFS đến Token-Gating (Bảng~\ref{tab:test_results}).

\begin{table}[htbp]
\centering
\caption{Kết quả kiểm thử tự động trên toàn bộ hệ thống}
\label{tab:test_results}
\begin{tabularx}{\textwidth}{l C C C}
\toprule
\textbf{Nhóm kiểm thử} & \textbf{Số lượng Kịch bản} & \textbf{Tỷ lệ Đạt} & \textbf{Thời gian Thực thi} \\
\midrule
Module Web3 (SIWE, IPFS, Token-Gate) & 4 & 100\% & 0.58 s \\
Module API Web \& Developer Portal & 41 & 100\% & 24.8 s \\
Module Lịch trình Đồng bộ (Sync Service) & 6 & 100\% & 4.2 s \\
Module Xử lý Dữ liệu Lịch & 107 & 100\% & 3.0 s \\
\midrule
\textbf{Toàn bộ hệ thống} & \textbf{158} & \textbf{100\% (158/158)} & \textbf{32.6 s} \\
\bottomrule
\end{tabularx}
\end{table}

Các kịch bản kiểm thử trọng tâm của kiến trúc Web3 bảo đảm:
\begin{itemize}
    \item Xác minh độ chính xác của thuật toán phục hồi địa chỉ ví từ chữ ký mật mã theo chuẩn EIP-4361 và khả năng nhận diện định dạng EIP-55.
    \item Kiểm tra tính toàn vẹn cấu trúc CIDv1 của dịch vụ IPFS trong cả chế độ ghim trực tiếp và chế độ giả lập xác định.
    \item Đánh giá tính chính xác của bộ lọc Token-Gate khi truy vấn số dư hợp đồng thông minh on-chain.
\end{itemize}

\subsection{Đánh giá độ sẵn sàng và hiệu năng dịch vụ}
\begin{itemize}
    \item \textbf{Độ sẵn sàng trong quá trình cập nhật:} Kiểm tra đo đạc thực nghiệm với chuỗi truy vấn liên tục vào endpoint kiểm tra sức khỏe trong suốt quá trình triển khai pod mới ghi nhận tỷ lệ phản hồi HTTP 200 đạt 100\%, chứng minh tính hiệu quả của chiến lược \texttt{RollingUpdate}.
    \item \textbf{Bảo mật danh tính và quyền riêng tư:} Hệ thống không lưu trữ bất kỳ thông tin nhận dạng cá nhân (PII) nào của người dùng. Khóa riêng tư ví được bảo vệ hoàn toàn ở phía client, loại bỏ rủi ro rò rỉ dữ liệu xác thực từ phía máy chủ.
\end{itemize}

\section{Kết Luận \& Hướng Nghiên Cứu Tiếp Theo}

\subsection{Kết luận}
Nghiên cứu đã hiện thực hóa thành công kiến trúc phi tập trung cho hệ thống quản lý và trích xuất lịch học tập đại học. Việc ứng dụng chuẩn định danh tự chủ EIP-4361 (SIWE), mạng lưu trữ phân tán IPFS, cơ chế kiểm soát truy cập bằng Smart Contract Token-Gating và hạ tầng vận hành Kubernetes Zero-Downtime đã chứng minh tính khả thi, độ tin cậy và tiềm năng của các công nghệ Web3 trong việc xây dựng các hệ thống phần mềm giáo dục tự chủ, minh bạch và có tính sẵn sàng cao.

\subsection{Hướng mở rộng trong tương lai}
\begin{enumerate}
    \item \textbf{Hỗ trợ giao thức WalletConnect v2:} Mở rộng khả năng tương thích với các ví Web3 trên thiết bị di động thông qua cơ chế quét mã QR an toàn.
    \item \textbf{Phát hành hợp đồng Soulbound Token (SBT):} Ứng dụng chuẩn NFT không thể chuyển nhượng (ERC-721 SBT) để tự động hóa quy trình phân quyền và hạn mức truy cập API dành cho sinh viên và nhà nghiên cứu.
    \item \textbf{Mã hóa phi tập trung với công nghệ Threshold Cryptography:} Ứng dụng các giao thức mã hóa ngưỡng phân tán (như Lit Protocol) nhằm đảm bảo chỉ có khóa riêng tư của người dùng mới có thể giải mã và đồng bộ dữ liệu phiên đào tạo của chính mình.
\end{enumerate}

\newpage
\bibliographystyle{plain}
\bibliography{references}

\end{document}